%*******************************************************
% 4.3.10 Độ tin cậy, giới hạn và khả năng thực hiện
%*******************************************************
\subsection{Độ tin cậy, giới hạn và khả năng thực hiện}\label{sec:tincay}

\begin{tomtat}
Với một sự kiện duy nhất, đề tài chỉ phát hiện được tác động đủ lớn: với các giả định minh họa, khoảng mức giảm 22\% của độ phi thanh khoản. Lực kiểm định phụ thuộc vào số tháng trước và sau, độ biến động và mức tự tương quan của chuỗi chênh lệch, gần như không phụ thuộc vào số cổ phiếu. Kết quả không ủng hộ giả thuyết vẫn được báo cáo đầy đủ. Năm phương án dự phòng được chuẩn bị cho các trục trặc dễ xảy ra nhất (Phụ lục~\ref{pl:ruiro}).
\end{tomtat}

\subsubsection*{Lực kiểm định}

Lực kiểm định của \HH{3} phụ thuộc chủ yếu vào số tháng trước và sau sự kiện, phương sai và mức tự tương quan của chuỗi chênh lệch theo tháng, và chỉ phụ thuộc yếu vào số cổ phiếu.

Điều này nghe ngược với trực giác, vì thông thường mẫu càng nhiều quan sát thì càng phát hiện được tác động nhỏ. Nhưng như mục~\ref{sec:kiemdinh} đã giải thích, toàn bộ thông tin được gộp thành một chuỗi theo tháng. Thêm cổ phiếu làm mỗi điểm của chuỗi bớt nhiễu riêng lẻ, nhưng không làm giảm các cú sốc chung của tháng, vốn chiếm phần lớn biến động của chuỗi. Thêm tháng thì có tác dụng thật.

Với giả định chuỗi chênh lệch theo tháng là một quá trình tự hồi quy bậc một có hệ số $\phi$ và độ lệch chuẩn không điều kiện $\sigma$, cửa sổ trước dài 193 tháng (6/2009--6/2025), 18 tháng chuyển tiếp bị loại và giai đoạn sau dài 24 tháng, hiệu ứng nhỏ nhất phát hiện được của $\delta$ ở lực kiểm định 80\% và mức ý nghĩa một phía 0,05/3 xấp xỉ $0{,}64\sigma$ khi $\phi = 0$, $1{,}24\sigma$ khi $\phi = 0{,}6$ và $1{,}76\sigma$ khi $\phi = 0{,}8$. Nếu cửa sổ trước được rút ngắn do điểm gãy cấu trúc, hiệu ứng nhỏ nhất phát hiện được tăng lên, khoảng 1,3 lần khi cửa sổ trước chỉ còn 24 tháng.

Chẳng hạn, với $\sigma = 0{,}2$ và $\phi = 0{,}6$, đề tài chỉ phát hiện được mức giảm logarit Amihud từ khoảng 0,25, tương đương độ phi thanh khoản giảm khoảng 22\%. Đây là giả định minh họa; $\sigma$, $\phi$ và hiệu ứng nhỏ nhất phát hiện được sẽ được ước lượng lại từ giai đoạn tham chiếu và ghi vào tệp đăng ký.

\begin{vidu}[: đọc bảng hiệu ứng nhỏ nhất phát hiện được]
Với $\sigma = 0{,}2$:
\begin{center}
\begin{tabular}{@{}cccc@{}}
\toprule
$\phi$ & Theo $\sigma$ & Theo logarit Amihud & Độ phi thanh khoản giảm\\
\midrule
0 & $0{,}64\sigma$ & 0,128 & khoảng 12\%\\
0,6 & $1{,}24\sigma$ & 0,248 & khoảng 22\%\\
0,8 & $1{,}76\sigma$ & 0,352 & khoảng 30\%\\
\bottomrule
\end{tabular}
\end{center}
Cột cuối tính bằng $1 - e^{-x}$ với $x$ ở cột thứ ba. Chuỗi càng \enquote{dính} giữa các tháng ($\phi$ lớn), mỗi tháng mới càng ít thông tin mới, nên tác động phải lớn hơn mới phát hiện được. Bảng dùng $\sigma = 0{,}2$ chỉ để minh họa.
\end{vidu}

\begin{delam}
\emph{Dễ nghĩ là} nếu không bác bỏ được giả thuyết gốc thì nghĩa là chuyển sàn không có tác động. Không đúng. Với lực kiểm định như trên, một tác động thật bằng 10\% rất có thể không được phát hiện. Kết luận đúng trong trường hợp ấy là \enquote{chưa đủ cơ sở kết luận}, kèm khoảng tin cậy cho biết những tác động nào còn phù hợp với dữ liệu.
\end{delam}

\HH{4} có ít cửa sổ nên chỉ được báo cáo ở dạng ước lượng, và so sánh theo hai giai đoạn căng thẳng chỉ ở dạng mô tả. Giá trị p đưa vào thủ tục Holm (1979) ở mức 5\% là giá trị p một phía: phân cụm theo quốc gia cho \HH{1}, bootstrap khối dừng cho \HH{2} và bootstrap tự hồi quy cho \HH{3}.

Ở \HH{1}, đề tài kiểm tra độ vững trên mẫu các trung tâm có ít nhất năm kỳ và trên mẫu cân bằng. Vì kỳ GFCI 40 đã có nhân tố trùng với $\SMD$ (Wardle \& Mainelli, 2026), \HH{1} được ước lượng lại trên các kỳ không trùng lặp nếu đủ số kỳ.

\subsubsection*{Khi dữ liệu không ủng hộ giả thuyết}

Nếu dữ liệu không ủng hộ giả thuyết, đề tài vẫn trả lời được câu hỏi nghiên cứu bằng kích thước hiệu ứng và khoảng tin cậy; khi khoảng tin cậy chứa cả 0 lẫn những tác động lớn hơn hiệu ứng nhỏ nhất phát hiện được, kết quả được gọi là chưa đủ cơ sở kết luận, không phải bằng chứng về việc không có tác động.

Có ba kiểu kết quả cần phân biệt. Khoảng tin cậy nằm hẳn dưới 0: tác động có, và có hướng như \HH{3} kỳ vọng. Khoảng tin cậy hẹp quanh 0, không chứa tác động nào lớn hơn hiệu ứng nhỏ nhất phát hiện được: có bằng chứng rằng tác động, nếu có, là nhỏ. Khoảng tin cậy rộng, chứa cả 0 lẫn tác động lớn: dữ liệu chưa đủ để nói. Kiểu thứ hai cũng là một phát hiện có giá trị chính sách.

\subsubsection*{Năm phương án dự phòng}

Kế hoạch có năm phương án dự phòng. Nếu không số hóa được hoặc không được phép dùng điểm GFCI theo kỳ, câu hỏi 1 được trả lời ở dạng mô tả và \HH{1} không được kiểm định. Nếu việc chuyển sàn bị lùi hoặc thực hiện theo nhiều đợt, đề tài dùng ngày thực tế của từng nhóm và bộ ước lượng khác biệt kép cho xử lý so le (Callaway \& Sant'Anna, 2021); nếu sự kiện không diễn ra, đề tài trả lời câu hỏi 3 bằng các mã đã tự chuyển từ HNX sang HOSE trước năm 2026.

Bộ ước lượng cho xử lý so le cần thiết vì khi các nhóm cổ phiếu chuyển sàn vào những ngày khác nhau, phương trình khác biệt kép thông thường có thể dùng nhóm đã chuyển trước làm đối chứng cho nhóm chuyển sau, và cho kết quả khó diễn giải. Callaway và Sant'Anna (2021) ước lượng tác động riêng cho từng nhóm theo ngày chuyển, chỉ so với các nhóm chưa chuyển, rồi mới gộp lại.

Thông báo của VNX chỉ áp dụng cho cổ phiếu niêm yết, nên đợt chuyển ngày 28/12/2026 không gồm UPCoM (Báo Đầu tư Chứng khoán, 2026; Sở Giao dịch Chứng khoán Việt Nam, 2026), nhưng hạn chuyển cổ phiếu UPCoM sang HOSE vẫn là hết ngày 31/12/2026 (Bộ Tài chính, 2023), nên UPCoM có thể chuyển sàn theo một đợt riêng trước hết ngày 31/12/2026, tức trước cửa sổ sau, hoặc trong cửa sổ sau nếu mốc này được lùi. Khi đó, nhóm đối chứng UPCoM và chỉ số của UPCoM chỉ được dùng đến tháng liền trước ngày chuyển, ước lượng chính dựa vào nhóm đối chứng HOSE, và tương tác theo biên độ dùng thước đo thay thế ở mục~\ref{sec:kiemdinh}.

Nếu cả hai nhóm đối chứng không đạt kiểm tra xu hướng song song, đề tài dùng nhóm đối chứng tổng hợp từ các mã HOSE ngoài danh sách FTSE Russell và mã UPCoM. Nếu dữ liệu sau sự kiện ngắn hơn 24 tháng hoặc rổ CS có ít hơn 30 cổ phiếu đủ thanh khoản, đề tài dùng, tương ứng, toàn bộ tháng có sẵn hoặc rổ 20 cổ phiếu.

Nhóm đối chứng tổng hợp là một tổ hợp có trọng số của nhiều cổ phiếu đối chứng, với trọng số được chọn sao cho tổ hợp đi sát nhất với nhóm chuyển sàn trong giai đoạn trước sự kiện. Nếu không có nhóm có sẵn nào đi song song, phương pháp này tạo ra một nhóm như vậy từ dữ liệu.
